\documentclass[preamble=tex/preamble.tex,authorsonlines]{markdown}
title: A sound pure type system with ranks as sorts and Type-in-Type
author:
- Pavel Sokolov:
    orcid: 0009-0001-6922-2870
    email: ppsokolov@hse.ru
    affiliation:
    - HSE University, Laboratory of Theoretical Computer Science
- Alexander Kuklev:
    orcid: 0009-0004-1019-4205
    email: a@kuklev.com
    affiliation:
    - Radboud Univ. Nijmegen, Department of Software Science
formalization:
  rocq-ai: https://github.com/TurtlePU/pts-rk/blob/8e3a7b67b571f6ab3fb7cc67e9dc7c0a214351ce/theories/
  rocq: https://github.com/TurtlePU/pts-rk/blob/d22053fa22d6070c8a14af4da679e439dd012825/theories/
msc2020: 03B38, 03D15, 03B70, 03B40, 68N18
date: October 16, 2026

:::Abstract
We present a dependently typed pure type system
that contains a circular axiom $s : s$ for a
certain sort $s$ and is nonetheless strongly
normalizing.
Girard's paradox is avoided because our sorts are
not closed under a $Π$-former:
function types are placed in sorts according to
their *rank*.
Turns out that this system has a unique set of
characteristics: it has similar normalization
behaviour as STLC (so complexity of normalization
is %textsf{TOWER}-complete) while at the same time
allowing non-uniform representation of arbitrary
elementary recursive functions; furthermore,
it has a nontrivial set-theoretical model in
ZF+GCH and allows to carry out nontrivial proof
developments.
Base normalization proof is surprisingly simple
(modulo general metatheoretical properties of pure
type systems) and has been fully formalized in
vanilla Rocq (formerly Coq) without additional
axioms.
:::

# Introduction

Pure type systems (PTS) are a family of minimalistic
typed $λ$-calculi introduced independently by
Berardi [@berardi1988towards] and Terlouw [@terlouw1989gstt].
Their abstract grammar includes only variables
$x ∈ \mathsf{Var}$, abstraction $λ(x : T).u$,
application $u\,t$, dependent function types
$Π(x : T).U$ (where $U$ might depend on $x$) and
*sorts* $s ∈ \mathcal{S}$ that act as “types of
types”.

:::Figure[ht] Grammar of pure type systems
%bnf
  $x$ : \textsf{Variable} :in: $X$;;
  $s$ : \textsf{Sort} :in: $S$;;
  $T$, $U$, $t$, $u$ : \textsf{Term} ::=
  | $x$ : variable
  | $s$ : sort as type
  | $Π (x \colon T).\,U$
      : type of dependent functions
  | $λ (x \colon T).\,u$ : anonymous function
  | $u\,t$ : function application
  ;;
:::

Common syntax conventions to make it more
lightweight are the following:
application associates to the left,
so $a\,b\,c$ is understood as $(a\,b)\,c$;
in a sequence of abstractions,
all $λ$ symbols except the first can be omitted;
$T → U$ is an alias for $Π (x : T).U$ (where $U$
does not depend on $x$); $( → )$ is
right-associative, so $T → U → V$ is understood as
$T → (U → V)$.
In addition, even though abstractions are always
annotated with a type of an argument,
we are going to omit it on paper whenever they are
clear from the context.
Also note that there is no syntactical distinction
between terms and types — they are defined by a
single grammar; an expression can act both as a
term and as a type depending on a side where it
stands next to the colon in the typing judgement
$Γ ⊢ \_ : \_$.

A single PTS is determined by two relations:

- *axiom* relation
  $\mathsf{Ax} ⊆ \mathcal{S} × \mathcal{S}$
  determines “to which sort(s) a given sort
  belongs as an element”; instead of
  $(s, s') ∈ \mathsf{Ax}$ it is customary to write
  simply $s : s'$.

- *rule* relation
  $\mathsf{Rule} ⊆ \mathcal{S}^2 × \mathcal{S}$
  places types of functions $Π (x : T).U$ into
  sorts accounting for the sorts where $T$ and $U$
  are situated.

:::Figure[h] Typing rules of pure type systems {#typing}
%mathpar
  \inferrule[Ax]{}{Γ ⊢ s : s'}
      (s, s') ∈ \mathsf{Ax}
  \and
  \inferrule[Rule]
    {Γ ⊢ T : s_T \\ Γ, x : T ⊢ U : s_U}
    {Γ ⊢ Π (x : T).\, U : s}
        (s_T, s_U, s) ∈ \mathsf{Rule}
  \and
  \inferrule[Var]{}{Γ ⊢ x : T}
      x : T ∈ Γ
  \and
  \inferrule[Abs]
    {Γ ⊢ T : s_T \\ Γ, x : T ⊢ u : U}
    {Γ ⊢ λ (x : T). u : Π(x : T). U}
  \and
  \inferrule[App]
    {Γ ⊢ u : Π(x : T). U \\ Γ ⊢ t : T}
    {Γ ⊢ u\,t : U[t/x]}
  \and
  \inferrule[Conv]
    {Γ ⊢ t : T \\ Γ ⊢ U : s \\ T =_β U}
    {Γ ⊢ t : U}
:::

Systems from Barendregt's lambda cube [@barendregt1991introduction; @barendregt1992lambda]
fit naturally into the framework: for example,
simply typed lambda calculus (STLC) [@church1940formulation]
can be viewed as a pure type system with a single
sort $∗$, no axioms and a single rule $(∗, ∗, ∗)$.
In general, however, pure type systems are
dependently typed since $Π$ allows the codomain
type to depend on the value of an argument —
Coquand's Calculus of Constructions (CoC) [@coquand1988calculus]
sits comfortably at the top of lambda cube and is
modelled as a pure type system with two sorts
$\{ ∗, □\}$, a single axiom $∗ : □$ and four rules
corresponding to simple functions $(∗, ∗, ∗)$,
polymorphic functions $(□, ∗, ∗)$,
type operators $(□, □, □)$ and dependent types
$(∗, □, □)$.

Both STLC and CoC are strongly normalizing and
sound with respect to their set-theoretical
semantics [@tait1967intensional; @coquand1990kripke; @geuvers1991modular],
but it is not true for arbitrary PTS.

:::theorem[Girard'72 \cite{girard1972interpretation, hurkens1995simplification}]
The following PTS called System U is not
    normalizing:

%mathpar
  \begin{array}{c}
    \mathcal{S} = \{ ∗, □, △\},
    \quad
    \mathsf{Ax} = \{(∗, □), (□, △)\},
    \\
    \mathsf{Rule} = \{
      (∗, ∗, ∗), (□, ∗, ∗), (△, ∗, ∗),
      (□, □, □), (△, □, □)
    \}
  \end{array}

Furthermore, the offending term has type
$Π (P : ∗).P$.
:::

This result is dubbed “Girard's paradox” since it
renders System U inconsistent with respect to its
intended logical semantics: inhabitedness of
$Π (P : ∗).P$ would mean that every proposition is
provable.
Thus, any system where System U embeds is
inconsistent, too; in particular,
an original version [@martinlof1971theory] of
Martin-Löf's Intuitionistic Type Theory [@martinlof1984intuitionistic]
embeds the following PTS known as “Type-in-Type”
which is inconsistent due to Girard's paradox:

%mathpar
  \mathcal{S} = \{ ∗ \}, \quad \mathsf{Ax} = \{ ∗ : ∗ \}, \quad
  \mathsf{Rule} = \{(∗, ∗, ∗)\}.

For a pure type system with usual set-theoretical
semantics, it is natural to expect that a sort of
*small types* $∗$ is closed under $Π$,
thus we cannot forbid the rule $(∗, ∗, ∗)$; then,
a circular axiom $∗ : ∗$ has to go.
Long story short, it motivates an introduction of
a *universe hierarchy* [@martinlof1975intuitionistic; @harper1991universes]
$\mathcal{S} = \{\mathcal{U}_k\}_{k ∈ ℕ}$ with
axioms $\mathcal{U}_k : \mathcal{U}_{k + 1}$ and
rules
$(\mathcal{U}_k, \mathcal{U}_l, \mathcal{U}_{k ∨ l})$
for each $k, l$; when interpreted in sets,
$⟦ \mathcal{U}_k ⟧$ is taken to be the $k$-th
inaccessible cardinal [@werner1997sets; @aczel1998relating].

Observing this, we took an interest in the following
two questions:

- If we were to drop the rule $(∗, ∗, ∗)$ and keep
  the axiom $∗ : ∗$ instead, can we find a
  nontrivial consistent pure type system?

- Is there a consistent pure type system with
  total relations $\mathsf{Ax}$ and
  $\mathsf{Rule}$ whose natural set-theoretical
  semantics do not involve the assumption of
  inaccessible cardinals?

Turns out that, for both questions,
the answer is “yes”.
Furthermore, both yes-s are witnessed by a single
PTS, which we dub %PTSrk{}:

%mathpar
  \mathcal{S} = \{\rk{r}\}_{r ∈ ℕ}
  \quad
  \mathsf{Ax} = \{(\rk{r}, \rk{0})\}_{r ∈ ℕ}
  \quad
  \mathsf{Rule} = \{(\rk{r}, \rk{s},
          \rk{(r + 1) ∨ s})\}_{r, s ∈ ℕ}

(Naming “$\mathsf{rk}$” comes from *rank*,
since rules compute precisely the rank of
functions, which is the maximal nesting of $Π$ in
the domain type.)

Indeed, there is a circular axiom
$\rk{0} : \rk{0}$, and System U is not embeddable
since $\rk{r}$ is not closed under $Π$ for any
$r$.
Also, as we will see in Section <#rk-sn>,
our %PTSrk{} is strongly normalizing despite an
intuition that $\rk{r} : \rk{0}$ might “feel
wrong”.
Modulo general metatheory of pure type systems
laid out in Section <#pts-metatheory>,
a normalization proof turned out to be
suspiciously simple, even more so in comparison
with normalization proofs of other dependently
typed systems [@geuvers1991modular; @barras1999autovalidation; @abel2018decidability].
To dispel our own doubts, we formalized a proof of
normalization in Rocq (formerly Coq) without using
any additional axioms, yielding a fully
constructive proof; our journey in formalization
is summarized in Section <#formalization> and the
whole development is available at <https://github.com/TurtlePU/pts-rk>.

Having completed the proof, we noticed that it
exposes a primitive recursive procedure to
normalize terms of %PTSrk{}; we dug further and
discovered that, actually, there is an intimate
connection between a class of elementary recursive
functions [@kalmar1943; @rose1984subrecursion] and
our system, as seen in Sections <#complexity>-<#computational>.
As a byproduct of our study on complexity,
we built a sound set-theoretic model for %PTSrk{}
(see Section <#rk-model>) which assumes
*Generalized Continuum Hypothesis* [@jech2003set; @hausdorff1908; @godel1938consistency; @godel1940consistency]
which is an unusual pick for dependently typed
systems where *inaccessible cardinals* are a more
common assumption [@werner1997sets; @lee2011proof].
Before concluding, we provide a sneak peek into
what it is like to write proofs in %PTSrk{} and
provide a full interpretation of
Tarski–Mostowski–Robinson arithmetic $R$ (Section <#proof-system>).

## Use of artificial intelligence

At various points in our research,
we interacted with GPT-6 Astra [@gpt6astra2026] to
prove something for us.
Sometimes it was a rehash of previously
established results; sometimes it presented
something new.
In case we used the ideas hinted by LLM,
we explicitly state so right next to the proof;
even so, the text is completely human written
since we wanted to make sure we completely
understood the results (in addition,
generated proofs always required additional
explanations or a complete rewrite to make them
understandable and simpler to read).c
We also used Claude Fable 5.1 [@claude_fable]
to identify additional points of reference.

## Formalization status

Every result up to the proof of strong
normalization for %PTSrk{} was formally verified
in Rocq using handwritten proofs.
Such proofs are marked with a %proof{rocq-ai:https://github.com/TurtlePU/pts-rk};
each icon is a hyperlink to the statement and
proof in Rocq with the version pinned at the moment of
submission.

Results after that were formalized in Rocq as
well, albeit not by us but rather by GPT-6 Astra.
Such results are marked with %formalized{rocq-ai:https://github.com/TurtlePU/pts-rk/pull/1};
just as with handwritten results, each icon is a
hyperlink to the statement and proof in Rocq with
the version pinned at the moment of submission.
Where only a component of a paper claim is
checked, a note explaining what exactly is
formalized accompanies the badge.

# Metatheory of pure type systems {#pts-metatheory}

For ease of presentation of our results in this
and following sections, we assume a Barendregt
convention [@barendregt1984lambda] that all bound
variables are distinct from free variables and
from each other; furthermore,
$α$-equivalence (renaming of bound variables) is
assumed as a trivial equality on terms.
Nevertheless, our formalization takes variable
treatment seriously; see Section <#formalization>
for more details.

## PTS as a rewriting system

In Fig.
<#typing>, a rule of *conversion* (%textsc{Conv})
mentions *definitional equality* of types
$T =_β U$.
It is defined as a reflexive, symmetric,
transitive closure of a one-step reduction.

:::Definition[One-step reduction]
Relation of a one-step reduction $t →_β t'$ is the
smallest relation satisfying the following rules:

%mathpar
  \inferrule{T →_β T'}
      {Π (x : T).U →_β Π (x : T').U}
  \and
  \inferrule{U →_β U'}
      {Π (x : T).U →_β Π (x : T).U'}
  \and
  \inferrule{T →_β T'}
      {λ (x : T).t →_β λ (x : T').t}
  \and
  \inferrule{t →_β t'}
      {λ (x : T).t →_β λ (x : T).t'}
  \and
  \inferrule{t →_β t'}{t\,u →_β t'\,u}
  \and
  \inferrule{u →_β u'}{t\,u →_β t\,u'}
  \and
  \inferrule[Beta]{}
      {(λ (x : T).t)\,u →_β t[u/x]}
:::

In essence, actual reduction happens only in rule
%textsc{Beta}; other rules simply state that we can
freely reduce any subterm.

The following result will be greatly useful in our
development.
It was proven many times over for different
similar systems [@church1936properties; @barendregt1992lambda];
we briefly revisit the proof since previous
presentations usually did not allow reductions in
the type of argument $T$ in $λ (x : T).t$.

:::Theorem[Characterization of $( =_\beta )$ \proof{rocq:DefEq.v\#L9}]
In an arbitrary PTS,

%mathpar
  t =_β u \quad ⇔ \quad
  ∃ v.\;t →_β^∗ v ∧ u →_β^∗ v,

where $( →_β^∗ )$ is the reflexive transitive closure
of $( →_β )$.
:::

:::Proof
Note that it is enough to prove that $( →_β^∗ )$
has a diamond property — for each $t_0, t_1, t_2$
such that $t_0 →_β^∗ t_1$ and $t_0 →_β^∗ t_2$
there exists a $t'$ such that $t_1 →_β^∗ t'$ and
$t_2 →_β^∗ t'$:
%vspace{-1em}

%tikzcd
  [ampersand replacement = \&]
    t_0 \arrow[r] \arrow[d] \&
    t_1 \arrow[d, dashed]
  \\
    t_2 \arrow[r, dashed] \&
    t'
Then, any ‘mountains’ in the path $t =_β u$ can be
switched into ‘valleys’.

Now, diamond property is proven by the standard
technique of parallel reduction [@martinlof1971theory]
which is defined as the smallest relation
satisfying:

%mathpar
  \inferrule{}{t ⇉ t}
  \and
  \inferrule
    {T ⇉ T' \\ U ⇉ U'}
    {Π (x : T).U ⇉ Π (x : T').U'}
  \and
  \inferrule
    {T ⇉ T' \\ t ⇉ t'}
    {λ (x : T).t ⇉ λ (x : T').t'}
  \and
  \inferrule
    {t ⇉ t' \\ u ⇉ u'}
    {t\,u ⇉ t'\,u'}
  \and
  \inferrule
    {t ⇉ t' \\ u ⇉ u'}
    {(λ (x : T).t)\,u ⇉ t'[u'/x]}

Intuitively, while $( →_β )$ is a single
$β$-reduction somewhere in the term and
$( →_β )^∗$ is a sequence of reductions,
$( ⇉ )$ is several simultaneous reductions
immediately available in the term.
Notice, then, that $( →_β ) ⊆ ( ⇉ ) ⊆ ( →_β^∗ )$;
by monotonicity of reflexive transitive closure,
$( →_β^∗ ) ⊆ ( ⇉^∗ ) ⊆ ( →_β^∗ )^∗ = ( →_β^∗ )$,
so $( ⇉^∗ ) = ( →_β^∗ )$.
Thus, it is enough to prove the diamond property of
$( ⇉^∗ )$.
Diamond property is preserved under reflexive
transitive closure, so it is enough to prove
diamond property of $( ⇉ )$, which holds by
tedious case analysis.
:::

Setting aside practical applications of this
theorem (to algorithmically check equivalence of
terms, we can now search for a common reduction
result), it will prove useful already in our proof
development.
For example, the following “equality inversion”
lemmas are direct corollaries:

:::Corollary[Equality inversion]
If $s =_β s'$, then $s = s' ∈ \mathcal{S}$;
if $Π (x : T).U =_β Π (x : T').U'$,
then $T =_β T'$ and $U =_β U'$.
:::

They will become necessary in proving subject
reduction, or type preservation.
But first, we will show some structural properties
of a typing judgement given in Fig.
<#typing>.

## Properties of a typing judgement

These properties have certainly been proven before
for PTS [@barendregt1992lambda; @geuvers1991modular],
but our formulation of the typing rules slightly
differs in important places, so it is better to
recheck that these properties hold.

:::Lemma[Weakening \proof{rocq:Typing.v\#L50}]
The following rule is admissible:

%mathpar
  \dfrac{Γ, Δ ⊢ t : T}
        {Γ, x : U, Δ ⊢ t : T}.
:::

:::Proof
By structural induction on the typing derivation.
:::

:::Lemma[Substitution \proof{rocq:Typing.v\#L78}]
The following rule is admissible:

%mathpar
  \dfrac{Γ, x : U, Δ ⊢ t : T \quad
         Γ ⊢ u : U}
        {Γ, Δ[u/x] ⊢ t[u/x] : T[u/x]}.
:::

:::Proof
By structural induction on the typing derivation
for $t$.
For variables in $Δ$, it is necessary that both
$Δ$ and $T$ get substitutions of $u$ for $x$ so
%textsc{Var} is still applicable;
for places in $t$ where $x$ used to be, we need to
apply a weakening lemma to the typing derivation for
$u$ exactly $|Δ|$ times.
In the case of %textsc{Conv} it is important that
definitional equality of types is preserved by
substitution (thanks to characterization of
$( =_β )$).
:::

In addition, there is a family of so-called
*inversion lemmas* [@barendregt1992lambda; @vanbenthemjutting1993typing]
stating that a type of well-typed term is largely
determined by the shape of a term.
They are also proven by straightforward induction
on typing derivation
(each one is proven separately,
no mutual induction involved):

:::Lemma[Inversion lemmas \proof{rocq:Inversions.v}] {#inversion}
%align
  Γ ⊢ s : T & ⇒ ∃ s'. (s, s') ∈ \mathsf{Ax} ∧ T =_β s'
  \\
  Γ ⊢ Π (x : T).U : V & ⇒ ∃ s_T, s_U, s.\left\{
  \begin{array}{l}
    (s_T, s_U, s) ∈ \mathsf{Rule} \\
    Γ ⊢ T : s_T \\
    Γ, x : T ⊢ U : s_U \\
    V =_β s
  \end{array}\right. \\
  Γ, x : T, Δ ⊢ x : U & ⇒ T =_β U
  \\
  Γ ⊢ λ (x : T).t : V & ⇒ ∃ s, U.\left\{
  \begin{array}{l}
    Γ ⊢ T : s \\
    Γ, x : T ⊢ t : U \\
    V =_β Π (x : T).U
  \end{array}\right. \\
  Γ ⊢ t\,u : V & ⇒ ∃ U, T.\left\{
  \begin{array}{l}
    Γ ⊢ t : Π (x : U).T \\
    Γ ⊢ u : U \\
    V =_β T[u/x]
  \end{array}\right.
:::

To progress further, we need to make certain
assumptions about the PTS in question.

:::Definition[Functional PTS]
A pure type system is called *functional* if
$\mathsf{Ax} ⊆ \mathcal{S} × \mathcal{S}$ and
$\mathsf{Rule} ⊆ \mathcal{S}^2 × \mathcal{S}$ are
functional relations, that is,

%mathpar
  (s : s') ∧ (s : s”) ⇒ s' = s”
  \quad \text{and} \quad
  (s_1, s_2, s) ∈ \mathsf{Rule} ∧ (s_1, s_2, s') ∈ \mathsf{Rule} ⇒ s = s'
:::

:::Definition[Total PTS]
A pure type system is called *total* if
$\mathsf{Ax} ⊆ \mathcal{S} × \mathcal{S}$ and
$\mathsf{Rule} ⊆ \mathcal{S}^2 × \mathcal{S}$ are
total relations, that is,

%mathpar
  ∀ s.∃ s'.(s, s') ∈ \mathsf{Ax}
  \quad \text{and} \quad
  ∀ s_1, s_2.∃ s.(s_1, s_2, s) ∈ \mathsf{Rule}
:::

:::Theorem[Uniqueness of typing \proof{rocq:Inversions.v\#L84}]
In a functional PTS, the following rule is
admissible:

%mathpar
  \dfrac{Γ ⊢ t : T \quad Γ ⊢ t : T'}
        {T =_β T'}.
:::

:::Proof
By structural induction on $t$ and inversion of
the given typing derivations.

- $t = s ∈ \mathcal{S}$: by functionality of
  $\mathsf{Ax}$.

- $t = Π (x : T).U$: inversion of typing gives a
  pair of (possibly different) sorts for $T$ and
  another pair for $U$.
  By induction we get that ranks inside pairs are
  definitionally equal, which by equality
  inversion gives strictly equal sorts.
  Then what is needed holds by the functionality of
  $\mathsf{Rule}$.

- The case for $t = x ∈ \mathsf{Var}$ is trivial.

- $t = (λ (x : T).u)$: by induction hypotheses and
  by congruence of $( =_β )$.

- $t = u\,v$: by typing inversion and by induction
  hypothesis for $u$ we get $u : Π (x : T_1).U_1$
  and $Π (x : T_1).U_1 =_β Π (x : T_2).U_2$.
  We need $U_1[v/x] =_β U_2[v/x]$ which we get
  from equality inversion and preservation of
  def.eq.
  under substitution. %vspace{-1em}
:::
%vspace{-.5em}
In our proof of normalization, we are going to
rely on yet another property: if there is a
derivation for $t : T$, then $T$ is itself
typeable with some sort.
However, it does not hold in our current
formulation of typing rules: in $Γ ⊢ x : T$,
nothing prevents a context $Γ$ from containing
an entry
$(x : T')$ where $T'$ is not typeable.
Thus, we need an additional assumption.

:::Definition[Context well-formedness]
Context $Γ$ is *well-formed* (written $⊢ Γ$) if
every type in it is typed by some sort.
Formally, a set of well-formed contexts is the least
set closed by the following rules:

%mathpar
  \dfrac{}{⊢ ⋅} \quad \text{and} \quad
  \dfrac{⊢ Γ \quad Γ ⊢ T : s}
        {⊢ Γ, x : T}

(where $⋅$ denotes an empty context).
:::

Usually, when presenting a dependently typed
system, context well-formedness is added as a
premise to typing rules which would otherwise be
premise-free (such as %textsc{Ax} and %textsc{Var}
in a PTS).
Here we instead provide it as an additional
assumption only when necessary;
in addition to leading to simpler metatheory (we
avoid an otherwise mutually inductive definition
of $( ⊢ \_)$ and $(\_ ⊢ \_ : \_)$),
this is closer to how decision procedures for
typing judgements are implemented.

:::Theorem[Typing yields correct type \proof{rocq:Preservation.v\#L65}] {#type-is-correct}
In a total PTS, if $Γ ⊢ t : T$ and $⊢ Γ$,
then there exists $s ∈ \mathcal{S}$ such that
$Γ ⊢ T : s$.
:::

:::Proof
By structural induction on the typing derivation.

- %textsc{Ax}, %textsc{Rule}: by totality of
  $\mathsf{Ax}$.

- %textsc{Var}: follows from $⊢ Γ$.

- %textsc{Abs}: by totality of $\mathsf{Rule}$.

- %textsc{App}: $Γ ⊢ t\,u : T[u/x]$.
  From induction hypothesis for $t$ we know that
  $Γ ⊢ t : Π (x : U).T : s$ for some $s$,
  then by typing inversion for $Π (x : U).T$ we
  get that $T$ is typable with some sort.
  Then what is needed follows from a substitution
  lemma.

- %textsc{Conv}: what is needed is present in the
  premise. %vspace{-1em}
:::
%vspace{-.5em}
Totality is not a strictly necessary requirement —
similar true propositions can be stated for
non-total PTS, for example, in CoC typing yields
type which is either itself typable or equal to
$□$ [@barendregt1992lambda; @luo1994computation].
But %PTSrk{} is indeed total so the theorem above
is enough for our purposes.

:::Theorem[Subject reduction \proof{rocq:Preservation.v\#L92}]
In a total PTS, the following rule is admissible:

%mathpar
  \dfrac{Γ ⊢ t : T \quad ⊢ Γ
         \quad t →_β t'}
        {Γ ⊢ t' : T}
:::

:::Proof
By structural induction over derivation of
$t →_β t'$.
Inversion of a typing judgement yields all
necessary info, albeit because of it types are now
given only up to definitional equality, so we have
to apply the conversion rule in each case.
Crucially, it requires the type we convert into to
be typable by a sort, so that is where we use
Theorem <#type-is-correct>.

The only tricky cases are reductions of a domain
type in $Π (x : T).U$ and $λ (x : T).t$.
Let us consider the case of the following
reduction:

%mathpar
  λ (x : T).t →_β λ (x : T').t.

To complete this case, we have to prove
that $Γ, x : T' ⊢ t : U$ given that
$Γ, x : T ⊢ t : U$.
By appropriately strengthening the statement,
we get that it is enough to prove admissibility of
the following rule called *context reduction*:

%mathpar
  \dfrac{Γ, x : U, Δ ⊢ t : T \quad
         ⊢ Γ, x : U, Δ \quad
         U →_β U'}
        {Γ, x : U', Δ ⊢ t : T}

The proof is, again, by induction on the typing
derivation.
The only non-trivial case is with %textsc{Var}
rule used on $x$; it is taken care of by
the appropriate use of %textsc{Conv} rule to convert
$x : U'$ to $x : U$.
To use conversion, we have to prove that $U$ is
well-typed; thankfully, it follows from
$⊢ Γ, x : U, Δ$.

Note that context reduction is admissible in
arbitrary PTS!
We do not need it anywhere else, though,
so we did not state it as a separate lemma.
:::

Characterization of definitional equality,
uniqueness of typing and subject reduction
together imply the following corollary we are
going to use a few times in our normalization
proof:

:::Corollary[Type stability]
In a total functional PTS, the following holds:

%mathpar
  \left\{
  \begin{array}{l}
    Γ ⊢ t : T\\
    Γ ⊢ t' : T'\\ ⊢ Γ \\
    t =_β t'
  \end{array}\right. \quad ⇒ \quad
      T =_β T'.
:::

## Properties of strongly normalizing terms

We are almost done with classical results!
What is left now is to state which terms are
strongly normalizing and collect properties of
s.n.
terms which allow our soundness proof to go
through.

Informally, a term is strongly normalizing iff
there are no infinite reduction chains starting
from it.
In constructive setting and in formalization,
the following definition is a bit better [@altenkirch1993formalization; @vanraamsdonk1999perpetual]:

:::Definition[Strongly normalizing term] {#sn-definition}
Set of s.n.
terms is the smallest set $\mathsf{sn}$ in which
the following rule is valid:

%mathpar
  \dfrac{∀ t', (t →_β t') ⇒ t' ∈ \mathsf{sn}}{t ∈ \mathsf{sn}}.

In other words, a term is strongly normalizing iff
all possible reductions from it are strongly
normalizing.
:::

Thus, terms which cannot be reduced further (also
known as terms *in normal form*) are strongly
normalizing; then, terms which directly reduce
into normal forms are also strongly normalizing;
and so on.
Therefore, by this definition,
strongly normalizing terms are those that every
reduction chain starting from them is finite.

The following properties are then provable by
definition of $\mathsf{sn}$ and by structural
induction on available derivations of
$\mathsf{sn}$ for terms in question:

:::Lemma[Structural properties of $\mathsf{sn}$ \proof{rocq:SN.v}] {#sn-props}
%mathpar
  \inferrule{}{x, s ∈ \mathsf{sn}}\and
  \inferrule
    {T ∈ \mathsf{sn} \\ U ∈ \mathsf{sn}}
    {Π (x : T).U ∈ \mathsf{sn}}\and
  \inferrule
    {T ∈ \mathsf{sn} \\ t ∈ \mathsf{sn}}
    {λ (x : T).t ∈ \mathsf{sn}}\and
  \inferrule{t\,u ∈ \mathsf{sn}}{t ∈ \mathsf{sn}}\and
  \inferrule{t[u/x] ∈ \mathsf{sn}}{t ∈ \mathsf{sn}}
:::

Application, of course, is the only nontrivial
case.
In the case of $β$-reduction, strong normalization of
both $(λ (x : T).t)$ and $u$ does not imply
anything about normalization behaviour of
$t[u/x]$: consider a classical example from
untyped lambda calculus [@barendregt1984lambda],
$ω = (λ (x : T).x\,x)$ and $Ω = ω \,ω$;
while not typable, $ω$ is strongly
normalizing, while $Ω$ decidedly is not.

The situation changes when the head of the application is
not an abstraction at all.
This leads us to the definition of an *atomic
term*.

:::Definition[Atomic term]
The set of atomic terms is the smallest set
that satisfies the following conditions:

1. $s$, $Π (x : S).T$ and $y$ are atomic for all
   variables $x$,$y$, sort $s$ and types $S$,$T$;

2. If $t$ is atomic, then $t\,u$ is also atomic.
:::

Atomic terms are useful precisely because their
strong normalization survives application:

:::Lemma {#sn-atomic}
If $t, u ∈ \mathsf{sn}$ and $t$ is atomic,
then $t\,u ∈ \mathsf{sn}$.%;%proof{rocq:SN.v\#L78}
:::

:::Proof
By induction on derivations of $\mathsf{sn}$ for
$t$ and $u$.
Consider possible redexes in $t\,u$:
$t\,u$ is not itself a redex since $t$ is atomic
and hence not a lambda; the only possible redexes
are either inside $t$ or inside $u$.
Applying corresponding induction hypotheses,
we reach the desired result.
The only caveat is that, for induction to go
through, we need to check that atomicity of $t$ is
preserved by reduction; this is indeed true by
induction on derivation of atomicity for $t$.
:::

A dual view on strong normalization of application
is the following: if we already know that
$t[u/x] ∈ \mathsf{sn}$, then the most tricky case
is already taken care of.

:::Lemma[Weak head expansion \cite{joachimski2003short, vanraamsdonk1995normalisation, altenkirch1993constructions}\;\proof{rocq:SN.v\#L167}] {#sn-whe}
%mathpar
  \left\{
  \begin{array}{l}
    t[u/x]u_1… u_n ∈ \mathsf{sn} \\
    U, u ∈ \mathsf{sn}
  \end{array}\right. ⇒ (λ (x : U).t)u u_1… u_n ∈ \mathsf{sn}.
:::

:::Proof
By induction on $n$.

- $n = 0$.
  We need to prove that if
  $t[u/x], u, U ∈ \mathsf{sn}$,
  then $(λ (x : U).t)u ∈ \mathsf{sn}$.
  Note that $t[u/x] ∈ \mathsf{sn}$ implies
  $t ∈ \mathsf{sn}$; then, by structural induction
  on derivations of $\mathsf{sn}$ for $t, u, U$ we
  get what we want by considering all possible
  reductions.

- $n → (n + 1)$.
  By the induction hypothesis, it is sufficient to show
  the following:

  \[
    \left\{
    \begin{array}{l}
      (λ (x : U).t)u u_1… u_n ∈ \mathsf{sn} \\
      t[u/x]u_1… u_{n + 1} ∈ \mathsf{sn}
    \end{array}\right. ⇒ (λ (x : U).t)u u_1… u_n u_{n + 1} ∈ \mathsf{sn}.
  \]
  This is true by structural induction on
  available evidence.%vspace{-1em}
:::

# Strong normalization of %PTSrk{} {#rk-sn}

:::Theorem[Strong normalization of \PTSrk{}\;\proof{rocq:Ranks.v\#L217}] {#normalization}
%mathpar
  Γ ⊢ t : T\, ∧ ⊢ Γ ⇒ t ∈ \mathsf{sn}.
:::

:::Proof
By induction on the typing derivation.
Case %textsc{Conv} is trivial; due to Lemma <#sn-props>,
cases %textsc{Ax}, %textsc{Rule},
%textsc{Var} and %textsc{Abs} present no problem
as well.
The only problematic case is application:

%mathpar
  \left\{
  \begin{array}{l}
    Γ ⊢ u : Π (x : T).U \\
    Γ ⊢ t : T \\ ⊢ Γ \\
    t, u ∈ \mathsf{sn}
  \end{array}\right. ⇒ u\,t ∈ \mathsf{sn}

Now, here is the part where stratification
according to ranks comes in useful.
As typing yields the correct type, there is some $r$
such that $Γ ⊢ Π (x : T).U : \rk{r}$.
Let us explicitly add this to the list of
premises; we will call the following proposition
*Application Lemma*:

%mathpar
  \mathrm{AL}_r ≜ \quad
         \text{“}\left\{
  \begin{array}{l}
    Γ ⊢ Π (x : T).U : \rk{r}\\
    Γ ⊢ u : Π (x : T).U \\
    Γ ⊢ t : T \\ ⊢ Γ \\
    t, u ∈ \mathsf{sn}
  \end{array}\right. ⇒ u\,t ∈ \mathsf{sn}\text{”}

What indexing by $r$ bought us is that now we can
try proving this lemma by induction on $r$.
Again, as it always is with strong normalization
of application, the problem is in case where
the application head is a lambda and a $β$-reduction
takes place.
This calls for another lemma,
*Substitution Lemma*:

%mathpar
  \mathrm{SL}_r ≜ \quad
         \text{“}\left\{
  \begin{array}{l}
    Γ, x : T, Δ ⊢ u : U\\
    Γ ⊢ T : \rk{r}\\
    Γ ⊢ t : T\\ ⊢ Γ \\
    t, u ∈ \mathsf{sn}
  \end{array}\right. ⇒ u[t/x] ∈ \mathsf{sn}\text{”}

A proof of $\mathrm{AL}_r$ for all $r$ (which is
what we are after) is achieved through complete
mutual induction with $\mathrm{SL}_r$.
Thus, there are two main parts:

- $(∀ s < r.\;\mathrm{SL}_s) ⇒ \mathrm{AL}_r$:
  by induction on derivations of
  $u, t ∈ \mathsf{sn}$.
  Consider possible reductions of $u\,t$:
  most of them happen either in $u$ or in $t$,
  so that results of reductions from $u\,t$ are
  strongly normalizing by inductive hypotheses.
  The only problematic case is when
  $u = λ (x : T').v$ and the reduction is
  $β$-reduction.
  By inversion of typing for $u$, it should be the
  case that $T' =_β T$; and by inversion of typing
  for $Π (x : T).U$, we get that $T : \rk{s}$ and
  $U : \rk{t}$ such that $r = \max(s + 1, t)$.
  Then, by type stability, $T' : \rk{s}$,
  and furthermore, $s < r$, so we can apply
  a substitution lemma.

- $(∀ s ≤ r.\;\mathrm{AL}_s) ⇒ \mathrm{SL}_r$.
  So we need to prove $u[t/x] ∈ \mathsf{sn}$.
  We do it by simultaneous induction on derivation
  of $u ∈ \mathsf{sn}$ and structural induction on
  $u$.
  As substitution propagates through $u$,
  most cases are satisfied by Lemma <#sn-props>.
  A funny case is when $u = x$: then,
  a premise that $t ∈ \mathsf{sn}$ helps.

  The only problematic case is
  application $u = u_1\,u_2$; we consider three
  cases here:

  - $u_1[t/x]$ is atomic — handled by Lemma <#sn-props>.

  - $u_1$ is not atomic — by virtue of it being
    well-typed, it means that
    $u = (λ (x : V).u')v_1 … v_k u_2$.
    All possible reductions from $u[t/x]$ are then
    handled by induction hypotheses.

  - $u_1$ is atomic, but $u_1[t/x]$ is not — by
    virtue of $u_1$ being well-typed,
    the only case when this is possible is when
    $u_1 = x v_1… v_k$, thus
    $u[t/x] = t (v_1[t/x])… (v_k[t/x])(u_2[t/x])$.
    From induction hypotheses we know that all
    $v_i[t/x]$ and $u_2[t/x]$ are s.n.;
    $t$ is s.n.
    by premise and its rank is $r$.
    Then $t v_1[t/x]$ is s.n.
    by application lemma for $r$;
    but note that rank of $t v_1[t/x]$ is not
    greater than $r$ so application lemma applies
    again.
    Thus, by induction on $k$, the whole chain of
    applications is s.n.%vspace{-1em}
:::
%vspace{-.5em}
This is a retelling of a proof that our friend
Aleksandr Levin got after half-jokingly entering
the following prompt into GPT-6 Pro:

> Consider the following pure type system:
> ```text
> S = { rk_n | n natural }
> Axiom = { (rk_n, rk_0) | n natural }
> Rule = { (rk_m, rk_n, rk_max(m+1,n)) | m, n natural }
> ```

> Is it strongly normalizing or not?

After approximately 20 minutes,
LLM outputted the proof which we retold above.
Looking at it in retrospect, it seems as if we
stratified functions by ranks with intent for such
a direct proof to go through; honestly,
we did not think of that and were eagerly trying
to adapt well-known approaches for proving
normalization of pure type systems employing
either saturated sets or candidates of
reducibility [@tait1967intensional; @girard1972interpretation; @girard1989proofs; @geuvers1991modular],
which did not work because our “hierarchy” of
ranks is not well-founded enough because of axioms
$\rk{n} : \rk{0}$.
This proof, however, does not even bother with it,
which means that we could pick *any* relation as
$\mathsf{Ax}$.
Nevertheless, putting all sorts into $\rk{0}$
seems the most sensible approach to us since,
informally, sort is kind of a constant which can
not be converted into a function.

## Formalizing normalization of %PTSrk{} in Rocq {#formalization}

Given the high technicality of the proof for
strong normalization which involves wrangling of
unusual definition for strong normalization and
complex induction patterns (and given
that proof was supplied by the probabilistic
machine), we decided to formalize it so that a
sound simple algorithm would check it for us.
One of the natural choices for formalization of
programming languages is Rocq [@rocq2025manual; @bertot2004interactive],
an interactive theorem prover based on
Polymorphic, Cumulative Calculus of Inductive
Constructions [@sozeau2014universe; @timany2018cumulative].
As is the case with similar typed calculi,
it is both a programming language and a
higher-order logic, and correctness of a proof is
checked by the type checker; in addition to that,
Rocq features powerful metaprogramming facilities
which allow writing custom *tactics* — programs
written in a meta-language which generate terms of
the target language, greatly simplifying the
proofs.
While set-theoretical semantics of Rocq assume the
existence of countable hierarchy of inaccessible
cardinals [@werner1997sets; @timany2017consistency],
our formalization never delves into the hierarchy
and can be formalized in a much more restricted
theory.

To formalize our strong normalization
proof, we first need to provide a general
metatheory of PTS as laid out in Section <#pts-metatheory>.
We could reuse existing projects on PTS
formalizations such as [@siles2012pure; @adams2006pure; @barras1997coqencoq; @felicissimo2024generic]
but we wanted to explore the topic ourselves, which
led to a few minor differences discussed
previously; thus we started from scratch.

When one formalizes syntax with binders in a proof
assistant, one of the most important decisions is
how to represent binders.
Here we went with making our terms indexed by the
number of free variables it uses,
a standard technique originating from [@bellegarde1994substitution; @bird1999debruijn; @altenkirch1999monadic]:

```coq
Inductive Term (L: Type) (n: nat): Type :=
| sort: L -> Term L n
| Π: Term L n -> Term L (S n) -> Term L n
| var: Fin n -> Term L n
| lam: Term L n -> Term L (S n) -> Term L n
| app: Term L n -> Term L n -> Term L n.
```

(Inside our formalization project,
this definition is located here: %proof{rocq:Syntax.v\#L7})

Here, $L$ is a type of sorts ($L$ comes from
*level* which is a name used for it in universe
hierarchies), and $n$ is a number of free
variables.
Entering a new scope (such as in $Π$ and in
$\mathtt{lam}$) adds one new free variable;
variables, then, have to be numbers in range
$[0;n)$, which are usually represented by a
$\mathtt{Fin}$ type.

What is convenient in this representation is that
it makes it harder to make mistakes when wrangling
with context, defining substitution etc. — all
the technical details which make working with
syntax with binders so hard (authors of the first
implementation of such calculi, AUTOMATH [@debruijn1972lambda; @debruijn1970automath],
remarked it being notoriously hard to get right).
Also, it has a certain category-theoretical feel to
it: sufficiently generalized operations of
renaming and substitution give the typal
representation of syntax a structure of a
*relative monad* [@altenkirch2015monads; @altenkirch1999monadic].
It is really hard to go wrong when implementing
these operations (and the natural implementation
turns variables into de Bruijn indices [@debruijn1972lambda]);
on top of them, the usual substitution is easily
definable.

Most of the proof development in Rocq mirrors what
is written in Section <#pts-metatheory>; in fact,
we first wrote a proof in Rocq and then retold it
in written English here, omitting tedious
technicalities.
The only considerable difference is in the handling of
the set of strongly normalizing terms since
the proof of Substitution lemma given
above does not follow the structure of inductive
definition <#sn-definition>.
Instead, we use a syntax-directed characterization
of the set of strongly normalizing terms [@Joachimski2003Short]:

:::Definition[Syntax-directed characterization of strong normalization \proof{rocq:SN.v\#L175}]
Sets of terms $\mathsf{SN}$ (strongly
normalizing), $\mathsf{SNx}$ (strongly normalizing
neutral) and a relation $→_\mathsf {WHR}$ (weak
head reduction) are the least sets where the
following rules are valid:

%mathpar
  \inferrule{}{\mathsf{SNx}\,s}(s ∈ \mathcal{S})
  \and
  \inferrule
    {\mathsf{SN}\,T \\ \mathsf{SN}\,U}
    {\mathsf{SNx}(Π (x : T).U)}
  \and
  \inferrule{}{\mathsf{SNx}\,x}
  \and
  \inferrule
    {\mathsf{SNx}\,t \\ \mathsf{SN}\,u}
    {\mathsf{SNx}\,(t\,u)}
  \\
  \inferrule{\mathsf{SNx}\,t}{\mathsf{SN}\,t}
  \and
  \inferrule
    {\mathsf{SN}\,T \\ \mathsf{SN}\,t}
    {\mathsf{SN}\,(λ (x : T).t)}
  \and
  \inferrule
    {t →_\mathsf {WHR}u \\ \mathsf{SN}\,u}
    {\mathsf{SN}\,t}
  \and
  \inferrule
    {\mathsf{SN}\,U \\ \mathsf{SN}\,u}
    {(λ (x : U).t)\,u →_\mathsf {WHR} t[u/x]}
  \and
  \inferrule{t →_\mathsf {WHR}t'}
            {t\,u →_\mathsf {WHR}t'\,u}
:::

In a sense, $\mathsf{SNx}$ is a description of
atomic strongly normalizing terms;
$\mathsf{SN}$ adds a case for lambdas and for weak
head expansion, covering all possible cases.
Existence of such sets is by virtue of them being
subsets of the set of terms, so such mutual
inductive definition is not problematic.

One can prove that set $\mathsf{SN}$ is equal to
$\mathsf{sn}$; we, however, need only the
following:

:::Lemma[Soundness of $\mathsf{SN}$ \proof{rocq:SN.v\#L222}]
$\mathsf{SN} ⊆ \mathsf{sn}$.
:::

:::Proof
By mutual induction on the evidence
$\mathsf{SN}\,t$ for arbitrary $t$.
To do this, we also have to provide a term-indexed
proposition we are proving for $\mathsf{SNx}$ and
a relation on terms we are proving for
$\mathsf{WHR}$:

- $∀ t, \mathsf{SNx}\,t ⇒ \mathsf{sn}\,t$.

- $∀ t\,t', \,(t →_\mathsf {WHR}t') ⇒ (t →_\mathsf {WHR - sn}t')$
  (where the last relation is defined in the same
  way as WHR, but the head contains evidence of
  $\mathsf{sn}$ instead of $\mathsf{SN}$.)

Proof is extremely simple: by Lemmas <#sn-props>,<#sn-atomic>,<#sn-whe>.
We just have to notice that if the term is
$\mathsf{SNx}$, it is atomic,
so lemma <#sn-atomic> applies in the relevant
case.
:::
Then, to prove strong normalization,
it is enough to prove that if a term is typable,
it is $\mathsf{SN}$; we only have to state
Application Lemma and Substitution Lemma in terms
of $\mathsf{SN}$, and the proof goes through.

:::Lemma
$\mathrm{AL}_s$ for all $s ≤ r$ implies
$\mathrm{SL}_r$. %proof{rocq:Ranks.v\#L74}
:::

:::Proof
We have to prove that $\mathsf{SN}\,t$ and
$\mathsf{SN}\,u$ imply $\mathsf{SN}\,(t[u/x])$.
We do so by mutual induction on evidence for $t$
with the following propositions proven for
$\mathsf{SNx}$ and WHR:

- $\mathsf{SNx}\,t ∧ Γ ⊢ t : T ⇒ \mathsf{SNx}\,(t[u/x]) ∨ (\mathsf{SN}\,(t[u/x]) ∧ ∃ s ≤ r. \;Γ ⊢ T : \rk{s})$;

- $(t →_\mathsf {WHR}t') ⇒ (t[u/x] →_\mathsf {WHR}t'[u/x]) ∧ (\mathsf{SN}\,(t'[u/x]) ⇒ \mathsf{SN}\,(t[u/x]))$.

The proof then follows the non-formal treatment
given above.
:::

# Computational complexity of %PTSrk{} {#complexity}

In this section, we prove that the convertibility
and the type checking of closed, well-typed
%PTSrk{} terms is
%textsc{TOWER}-complete.

(Proofs in this and following sections were not
verified in Rocq by any of the authors,
but there is a formalization of some results by
GPT-6 Astra. Informal proofs, however,
are relatively simple, non-technical, and their
formalization requires many times more preliminary
work than seems reasonable.)

As requested further, LLM confessed that it was
largely inspired by studies on the complexity of
normalization for STLC [@statman1979typed; @schwichtenberg1991upper; @beckmann2001exact; @nguyen2024simply].
And indeed, upon taking a closer look at the proof
of Theorem <#normalization>, one cannot help but
notice it actually being a primitive recursive
algorithm [@skolem1923; @kleene1952introduction; @peter1951rekursive]
for term normalization in disguise.
Or something a tad faster:

:::Definition[$d$-bounded terms]
A %PTSrk{} term is $d$-bounded if it is well-typed
in some well-formed context, and the heads of all
redexes in it have rank equal to or below $d$.
:::

:::Theorem[Elementary recursive normalization of bounded \PTSrk{} \formalized{rocq-ai:NormalizationComplexity/BoundedNormalization.v\#L19}] {#bounded-sn}
For each $d$, there is an elementary recursive
algorithm [@rose1984subrecursion] for finding a
normal form of $d$-bounded terms.
:::

:::Proof
Consider the following relation of “greedy
parallel reduction”:

%mathpar
  \inferrule{}{x ⇒ x\\ s ⇒ s}\and
  \inferrule
    {T ⇒ T' \\ U ⇒ U'}
    {Π (x : T).U ⇒ Π (x : T').U'}\and
  \inferrule
    {T ⇒ T' \\ t ⇒ t'}
    {λ (x : T).t ⇒ λ (x : T').t'}\and
  \inferrule
    {t ⇒ (λ (x : T).t') \\ u ⇒ u'}
    {t\,u ⇒ t'[u'/x]}\and
  \inferrule
    {t ⇒ t' \\ t' ≠ λ … \\ u ⇒ u'}
    {t\,u ⇒ t'\,u'}

It differs from the usual parallel reduction since it
greedily applies every “many-argument” redex
available in a term.
For example,
$(λ (x : X).λ (y : Y).λ (z : Z).t)\,a\,b\,c$ gets
reduced to $t'[a/x][b/y][c/z]$ in a single step.
This relation has the following properties:

- $( ⇒ ) ⊆ ( →_β^∗ )$, so it is several usual
  reductions baked into one.

- This relation is total and functional.

- This relation can be realized with an elementary
  recursive algorithm $\mathcal{A}$:
  traverse the term bottom-up, immediately
  applying every $β$-redex which is formed when
  the result of a recursive run of the algorithm on
  the application head is a lambda.

- If $t ⇒ t'$, then
  $|t'| ≤ (2 + ε )^{\mathcal{O}(|t|)}$.

- If redex heads in $t$ have ranks less than *or
  equal to* $d$ and $t ⇒ t'$, redex heads in $t'$
  have ranks *less than* $d$: new redexes appear
  only after substitution of abstractions into
  application heads, and ranks of these
  abstractions are strictly less than the rank of
  redex heads in which these substitutions
  happened.

Because of the last property, it is enough to run
$\mathcal{A}$ for $d$ times in total to normalize
a $d$-bounded term; since input size on $k$-th run
is bounded by $k$-th tower of exponents,
the whole algorithm is elementary recursive.

To make sure that the corresponding function is
total, output $0$ if the result of an algorithm
contains a $β$-redex.
:::

Similar result was already proven for STLC [@schwichtenberg1991upper; @beckmann2001exact]
since simple types can be classified by ranks as
well.
And since rank is easily computable for simple
types, normalization of arbitrary STLC terms is in
$\mathsf{TOWER}$.
A little weaker result holds for our dependently
typed PTS:

:::Theorem {#tower-convertibility}
Convertibility of **closed** well-typed %PTSrk{}
terms is in $\mathsf{TOWER}$.
:::

:::Proof
By Theorem <#bounded-sn>, it is enough to compute
some upper bound $B$ on the rank of types
appearing in abstractions of a term,
then run a normalizer for $(B + 1)$-bounded terms
and compare the results; $B$ has to be bound by an
elementary recursive function on the size of
input. We present such a function in
Figure <#rank-bound-algorithm>:
$\mathtt{estimateRank}\;t$ provides a bound on
ranks of redex arguments in $t$.

:::Figure[htb] Computing a bound on ranks of redex arguments {#rank-bound-algorithm}
```agda
maxRankLiteral rkᵣ      = r
maxRankLiteral Π(x:T).U =
    maxRankLiteral T ⊔ maxRankLiteral U
maxRankLiteral x        = 0
maxRankLiteral λ(x:T).t =
    maxRankLiteral T ⊔ maxRankLiteral t
maxRankLiteral (t u)    =
    maxRankLiteral t ⊔ maxRankLiteral u

traversal _ rkᵣ      _ = (r , 0 , 0)
traversal Γ Π(x:T).U d =
    let (vᵀ , rᵀ , aᵀ) = traversal Γ T d
        (vᵁ , rᵁ , aᵁ) = traversal (x := vᵀ, Γ) U d
     in (vᵁ , suc rᵀ ⊔ rᵁ , aᵀ ⊔ aᵁ)
traversal Γ x        d = (d , Γ[x] , 0)
traversal Γ λ(x:T).t d =
    let (vᵀ , rᵀ , aᵀ) = traversal Γ T d
        (vᵗ , rᵗ , aᵗ) = traversal (x := vᵀ, Γ) t d
     in (vᵗ , rᵗ , rᵀ ⊔ aᵀ ⊔ aᵗ)
traversal Γ (t u)    d =
    let (vᵗ , rᵗ , aᵗ) = traversal Γ t d
        (_  , _  , aᵘ) = traversal Γ u d
     in (vᵗ , rᵗ , aᵗ ⊔ aᵘ)

estimateRank t =
    let (vᵗ , rᵗ , aᵗ) =
        traversal [] t (maxRankLiteral t)
     in aᵗ
```
:::
The function $\mathtt{maxRankLiteral}\;t$ computes
maximal rank which is contained as a subterm in
the term $t$.
The workhorse of this algorithm,
$\mathtt{traversal}\;Γ \;t\;d$,
accepts the list $Γ$ of ranks in context of a term
$t$ and a “default rank” $d$,
and outputs the triple $(v, r, a)$,
where $v$ is value of $t$ as a rank (if supplied
with enough arguments), $r$ is its rank if $t$ is
a type (if supplied with enough arguments) and $a$
is an upper bound on rank of types appearing in
abstractions which we are trying to compute.
The only suspicious case is
$\mathtt{traversal}\,Γ \,x\,d$ where we output $d$
(the “default” rank) as a value of $x$:
could we compute the value of $x$ instead,
if present?
Regrettably, this would require normalization of a
term corresponding to $x$, which is too long;
instead we notice that $\rk{r}$ can be reduced
only from terms which already contain it as a
subterm, so the value of $x$ as a rank can only be
$\rk{r}$ which is already present somewhere in
the original term.
And our result is monotone with respect to the
value we output here, so it is enough to provide
some upper bound on it; thus it is enough to
output a maximal rank which occurs in the original
term as a subterm.
:::
Closedness is essential: without the context,
we cannot determine the ranks of free type
variables and hence cannot bound
redex ranks.

:::Theorem
Convertibility of closed well-typed %PTSrk{} terms
is %textsc{TOWER}-hard (hence
%textsc{TOWER}-complete). %formalized{rocq-ai:NormalizationComplexity/STLCEmbedding.v\#L198}
^[The badge certifies the annotation-preserving
reduction from closed Church-style STLC,
including typing, preservation and reflection of
conversion, and the exact additive encoding-size
bound.]
:::

:::Proof
By reduction from convertibility of Curry-style
STLC terms, since that problem is
%textsc{TOWER}-hard [@nguyen2024simply].

Just in case, we start from Curry-style STLC
terms.
We then reduce this task to convertibility of
**closed** Curry-style STLC terms by adding
$λ$-abstractions to close over free variables;
finally, we infer types of variables in an STLC
term by a linear unification algorithm referenced
in [@nguyen2024simply] and settle on **closed**
well-typed **Church-style** STLC terms.
The following translation then yields **open**
well-typed %PTSrk{} terms:

%align*
  ⦇ x ⦈ & = x \\
          ⦇ λ x^σ.t
          ⦈ & = λ (x : σ).
              ⦇ t ⦈ \\
          ⦇ t\,u ⦈ & = ⦇ t ⦈ \,
          ⦇ u ⦈

The problem is that types $σ$ might contain type
variables $T_1, …, T_n$ which are not bound.
However, it is easily fixable by introducing
missing types via abstractions:

%mathpar
  t ↝ λ (T_1 : \rk{0})…
     (T_n : \rk{0}).⦇ t ⦈

Assigning all type variables rank $\rk{0}$
preserves well-typedness since values of such
types were never treated as functions in the original
term.
:::

:::Corollary
Type checking for closed %PTSrk{} terms is
%textsc{TOWER}-complete.
:::

:::Proof
The following term demonstrates that
convertibility of closed well-typed terms
$t, t' : T$ reduces to type checking:

%mathpar
  λ (P : T → \rk{0})(x : P\,t).
  (λ (y : P\,t').y)\,x

This term is well-typed iff $P\,t =_β P_t'$ which
holds iff $t =_β t'$.

Therefore, we only need to prove that type checking
for %PTSrk{} is in %textsc{TOWER}.
Consider the following syntax-directed type
synthesis rules:

%mathpar
  \inferrule[Ax']{}
  {Γ ⊢ \rk{r} ⇒ \rk{0}}
  \and
  \inferrule[Rule']
    {Γ ⊢ T ⇒ \rk{r} \\ Γ, x : T ⊢ U ⇒ \rk{s}}
    {Γ ⊢ Π (x : T).\,U ⇒ \rk{\max(r + 1, s)}}
  \\
  \inferrule[Var']{}
  {Γ ⊢ x ⇒ T}x : T ∈ Γ
  \and
  \inferrule[Abs']
    {Γ ⊢ T ⇒ \rk{r} \\ Γ, x : T ⊢ u ⇒ U}
    {Γ ⊢ λ (x : T).u ⇒ Π (x : T).U}
  \and
  \inferrule[App']
    {Γ ⊢ u ⇒ Π (x : T).U \\ Γ ⊢ t ⇒ T' \\ T =_β T'}
    {Γ ⊢ u\,t ⇒ U[t/x]}

This relation is functional (in the sense that if
$Γ ⊢ t : T$ and $Γ ⊢ t : T'$,
then $T = T'$ as words), and the corresponding
partial function is realized by an algorithm in
%textsc{TOWER}: this algorithm recurses on the
term $t$, matches on the inferred types as needed
and collects equalities on types it has to check.
Then, it can naïvely check all equalities in bulk:
there are at most $\mathcal{O}(|t|)$ such
equalities on terms of size
$\mathcal{O}(2^{\mathrm{poly}(|t|)})$,
so this runs in %textsc{TOWER} by Theorem <#tower-convertibility>.
To finish the proof, it is enough to notice that

%mathpar
  Γ ⊢ t : T \quad ⟺ \quad
  ∃ T', r.\;(Γ ⊢ T : r) ∧ (T =_β T') ∧ (Γ ⊢ t ⇒ T').

So the main algorithm infers the type of $T$,
checks that it is $\rk{r}$ for some $r$,
infers the type of $t$ and compares it with $T$.
:::

# %PTSrk{} as a computational model {#computational}

In this section we prove that functions
representable in %PTSrk{} are,
under minor assumptions, precisely elementary
recursive functions, and provide a bunch of other,
more specialized, results.

Another consequence of Theorem <#bounded-sn> is a
natural upper bound for functions *uniformly*
representable in %PTSrk{} (this result is also
similar to the one for STLC [@schwichtenberg1976definierbare; @statman1979typed; @fortune1983expressiveness]).

:::Definition
A family of functions on natural numbers
$\mathcal{F}$ is *uniformly* representable in
%PTSrk{} if there is an *input natural number
type* $N_I : \rk{r}$ and an *output natural number
type* $N_O : \rk{s}$ together with *encodings*
$∀ n ∈ ℕ.[n]^I : N_I$ and $∀ n.[n]^O : N_O$ such
that all of them are in normal form and for each
$(f : ℕ^k → ℕ) ∈ \mathcal{F}$ there is a
*representation* $F : N_I → … → N_I → N_O$ which
is also in normal form and such that

%mathpar
  ∀ (n_1, …, n_k) ∈ ℕ^k : \quad F\,[n_1]^I\,…\,[n_k]^I =_β [f(n_1, \,…, \,n_k)]^O
:::

:::Remark
Note that the requirement for terms to be in normal
form is not restrictive: we can always normalize a
non-normal term.
:::

:::Theorem {#uniform-elementary}
A family of elementary recursive functions is
**not** uniformly representable in %PTSrk{}.
:::

:::Proof
By a direct diagonalization argument.
Suppose we could encode all elementary recursive
functions with $N_I : \rk{r}$ and $N_O : \rk{s}$.

By Theorem <#bounded-sn>, the normalization of
$(r + 1)$-bounded terms $\mathcal{SN}_{r + 1}$ is
elementary recursive.
Therefore, the following function is also
elementary recursive:

%mathpar
  h(t) =
  \begin{cases}
    1, & \mathcal{SN}_{r + 1}(t\,[\mathtt{t}]^I) = [0]^O, \\
    0 & \text{otherwise}.
  \end{cases}

(Where $\mathtt{t} ∈ ℕ$ is code of a term $t$ used
in the definition of $h$.)

By assumption, it is encodable by $H : N_I → N_O$.
Observe the value of $h(H)$:

%mathpar
  h(H) = 1 ⇔ \mathcal{SN}_{r + 1}(H\,[\mathtt{H}]^I) = [0]^O ⇔ H\,[\mathtt{H}]^I =_β [0]^O ⇔ h(H) = 0

The first equivalence is by definition of $h$;
the second is by virtue of expression under
$\mathcal{SN}$ being $(r + 1)$-bounded;
the last is by definition of $h$ being encoded by
$H$.
:::

The next logical question is:
which functions are computable?
Elementary recursive!
Though, of course, not in the same sense as
before.

:::Definition[Polymorphic numeral family]
For each fixed $j ≥ 0$, there is a type of
*polymorphic Church numerals* of rank $j + 2$:

%mathpar
  ℕ_j ≜ Π (X : \rk{j}).
         (X → X) → X → X;

and for each fixed $n$, there is a corresponding
*Church literal* of type $ℕ_j$:

%mathpar
  \overline{n} ≜ λ \,X\,f\,x.f^n\,x.

($f^n\,x$ means repeated application of $f$ to
$x$: $f(…(f\,x)…)$.)
:::

:::Theorem[Non-uniform elementary representability] {#non-uniform-rep}
For each elementary recursive function
$f : ℕ^n → ℕ$, there is a *cost* $c_f$ such that
for each *base rank* $j$, there is a
*representation*
$F_j : ℕ_{j + c_f} → … → ℕ_{j + c_f} → ℕ_j$
such that

%mathpar
  ∀ (x_1, …, x_n) ∈ ℕ^n : \quad
  F_j\,\overline{x_1}\,…\,\overline{x_n} =_β \overline{f(x_1, …, x_n)}.
:::

:::Proof
As is well-known, simply typed Church numerals
admit addition and multiplication [@church1941calculi; @barendregt1984lambda].
It is still true for our rank-refined version:

%mathpar
  \begin{array}{rcll}
    \mathsf{Add}_j & ≜ &
    λ \,m\,n\,X\,f\,x.
    m\,X\,f\,(n\,X\,f\,x) & : ℕ_j → ℕ_j → ℕ_j \\
    \mathsf{Mul}_j & ≜ &
    λ \,m\,n\,X\,f\,x.
    m\,X\,(n\,X\,f)\,x & : ℕ_j → ℕ_j → ℕ_j
  \end{array}

By substitution-basis theorem [@prunescu2026minimal],
it is enough to additionally provide modulus
$x\,\mathrm{mod}\,y$ (where
$x\,\mathrm{mod}\,0 = x$) and binary
exponentiation $2^x$ since substitution,
constants and projections are “free” in our
calculus.

By Theorem <#uniform-elementary>,
we cannot hope to have both remaining operations
represented uniformly.
Indeed, naïve import of traditional exponentiation
term for untyped Church numerals [@church1941calculi; @rosser1984highlights]
raises the rank of an exponent:

%multline*
  \mathsf{Exp}_j ≜ λ \,m\,n\,(X : \rk{j})\,f\,x.\,
          n\,(X → X)\,(λ (h : X → X)(y : X).
              m\,X\,h\,y
          )\,f\,x
          \\
          : ℕ_j → ℕ_{j + 1} → ℕ_j

And that is exactly the reason we allowed input
and output ranks of Church numerals to differ.

There still remains a problem, however:
by allowing ranks to differ, we lost uniformity of
substitution.
In particular, consider an elementary function
$x ↦ x + 2^x$.
To represent such computation using combinators we
defined, $x$ has to have both type $ℕ_j$ (as
left addend) and $ℕ_{j + 1}$ (as
exponentiation power).
Instead of introducing subtyping in the system,
we make use of the following family of
*numeral-preserving lowerings*:

%mathpar
  δ_j ≜ λ \,n\,X\,f\,x.\,
          n\,(X → X)\,(λ \,h\,y.f\,(h\,y))\,
          (λ \,y.y)\,x : ℕ_{j + 1} → ℕ_j

As the name implies, the goal of such combinators
is to lower the rank of a numeral,
but preserve its “value”.
It is achieved in the following way:
given an $f$ and $x$ for lower-ranked type,
an original combinator for $n$ is used to build an
$n$-fold composition of $f$ with itself, which is
then applied to $x$.

Thus, we can lower ranks arbitrarily:
$δ_{b, a} ≜ δ_a ∘ … ∘ δ_{b - 1} : ℕ_b → ℕ_a$.
This restores substitution and allows gluing
representation of elementary recursive functions
together; cost $c_f$ of a function can be computed
from any expression that represents it — addition
and multiplication cost $0$, exponentiation costs
$1$, operands are merged by taking the maximum of
their costs.
(Of course, cost is not unique and depends on the
way the function is represented with an expression.)

All that is left, then, is to represent modulus.
The general idea is as follows: first,
suppose that we are dealing with a simpler problem
of taking modulus by a constant modulo,
e.g. $\mathrm{mod}\,5$.
Notice that there are only $5$ possible answers:
$0$,$1$,$2$,$3$,$4$.
As such, they can be represented by the type which
contains exactly $5$ elements;
Church encoding [@bohm1985automatic] yields a type
of *selectors*:

%mathpar
  V_5 ≜ Π (X : \rk{j}).\,
          X → X → X → X → X → X : \rk{j + 1}

Its $5$ different canonical values simply pick one
of $5$ inputs, to each their own.

There is an operation of *cyclic shift* on
selectors:

%mathpar
  c_5 ≜ λ \,s\,X\,x_1\,x_2\,x_3\,x_4
      \,x_5.\,s\,X\,x_2\,x_3\,x_4\,x_5\,x_1 : V_5 → V_5

Iterating it $n$ times on some pinned “zeroth”
selector, we actually get the
$(n\,\mathrm{mod}\,5)$-th selector.
All that is left is to read the numeral back:

%mathpar
  ρ ≜ λ \,s\,X\,f\,x.\,s\,X\,
      (f^0\,x)\,(f^1\,x)\,(f^2\,x)\,(f^3\,x)\,(f^4\,x) : V_5 → ℕ_j

Now, to model $m\,\mathrm{mod}\,n$ in
this vein, we need to *generate* $V_n$
dynamically.
Lucky for us, %PTSrk{} is dependently typed so
there is hope.

Indeed, the following construction was
produced in a discussion with GPT-6 Astra.

We are going to produce a combinator
$\mathsf{Mod} : ℕ_{j + 1} → ℕ_{j + 4} → ℕ_j$.
Fix parameter $X : \rk{j}$ and consider the
following signature for abstract $S : \rk{j + 1}$:

%mathpar
  \begin{array}{ll}
    K : X → S, & z : S, \\
    L : (X → S) → X → S, & c : S → S, \\
    ρ : S → (X → X) → X → X,
    & r : (X → X) → (X → X).
  \end{array}

Intention is that $S$ is simply-typed $V_q$ for
$X$, $z$ is its “zero”, $K$ is a “zero” for
$V_{q + 1}$, $c$ is a cyclic shift for $V_q$,
$L$ is a cyclic shift for $V_{q + 1}$,
$ρ$ converts a selector into a numeral,
and $r$ is the result of modulus computed using
structure from the previous step (this is similar
to the trick used to define a predecessor
function).

A callback accepting such a structure and returning
$X$ has the following type:

%align*
  H_X ≜ \; & Π (S : \rk{j + 1}).\;
          (X → S) → S → ((X → S) → X → S) → (S → S) → \\ & (S → (X → X) → X → X) → ((X → X) → X → X) → X
          \quad : \rk{j + 3}

A type of “records” is thus
$P_X ≜ H_X → X : \rk{j + 4}$,
with a constructor

%mathpar
  ⟨ S, K, z, L, c, ρ, r⟩ ≜ λ h.\,h\,S\,K\,z\,L\,c\,ρ \,r.

As an initial $P_X$, we build a structure for
$V_1$ and an $r$ which behaves like
$m\,\mathrm{mod}\,0 = m$:

%align*
  S_0 & ≜ X → X, \\
          K_0 & ≜ λ (a : X)(b : X).a, &
          z_0 & ≜ λ (a : X).a, \\
          L_0 & ≜ λ (t : X → S_0)(a : X)(b : X).
              t\,b\,a, &
          c_0 & ≜ λ (s : S_0).s, \\
          ρ_0 & ≜ λ (s : S_0)(g : X → X)(a : X).s\,a, &
          r_0 & ≜ δ_j\,m\,X.

Pack them into $p_0 ≜ ⟨…⟩$.

Then, given a structure for $q$,
the structure for $q + 1$ is as follows:

%align*
  S^+ & ≜ X → S, \\
          K^+ & ≜ λ (a : X)(b : X).Ka, &
          z^+ & ≜ K, \\
          L^+ & ≜ λ (t : X → S^+)(a : X)(b : X).
              L\,(t\,b)\,a, &
          c^+ & ≜ L, \\
          ρ^+ & ≜ λ (t : S^+)(g : X → X)(a : X).
              ρ \,(ta)\,g\,(ga), &
          r^+ & ≜ ρ \,(m\,S\,c\,z).

A transition term can be written like this:

%mathpar
  T ≜ λ (p : P_X)(h : H_X).
      p\,(λ \,S\,K\,z\,L\,c\,ρ \,r.\,
      h\,S^+\,K^+\,z^+\,L^+\,c^+\,ρ^+\,r^+) : P_X → P_X

Therefore, the following term computes the modulus
with the expected behaviour.

%multline*
  \mathsf{Mod}_j ≜ λ
          (m : ℕ_{j + 1})(n : ℕ_{j + 4})
          (X : \rk{j})(f : X → X)(x : X).\\
              n\,P_X\,T\,p_0\,
              (λ \,\_\,\_\,\_\,\_\,\_\,\_\,r.
                  r\,X\,f\,x).
:::

:::Remark
Note that the only difference between uniform and
non-uniform representations is that,
in non-uniform representations, we are allowed to
pick different $N_I$ and $N_O$ for each function.
:::

:::Corollary
Unary functions in $\mathcal{E}^2$ are uniformly
representable in %PTSrk{}.
:::

:::Proof
There is a binary function $\mathsf{ev} : ℕ^2 → ℕ$
in $\mathcal{E}^3$ which is universal for unary
functions in $\mathcal{E}^2$ [@Skordev2008, Section 5, p. 873, footnote 5].
In other words, for each $f : ℕ → ℕ$ in
$\mathcal{E}^2$, there is a code $\mathtt{f} ∈ ℕ$
such that

%mathpar
  ∀ n, \quad \mathsf{ev}(\mathtt{f}, n) = f(n).

And $\mathcal{E}^3$ is precisely elementary
recursive functions, so by Theorem <#non-uniform-rep>
there is a cost $r$ and a term
$E : ℕ_r → ℕ_r → ℕ_0$ such that

%mathpar
  ∀ m n, \quad E\,\overline{m}\,\overline{n} =_β \overline{\mathsf{ev}(m, n)}

Therefore, there exist $N_I ≔ ℕ_r$ and
$N_O ≔ ℕ_0$ such that, for each unary
$f ∈ \mathcal{E}^2$, there is a term
$F ≔ E\,\overline{\mathtt{f}}$ such that

%mathpar
  ∀ n, \quad F\,\overline{n} =_β E\,\overline{\mathtt{f}}\,\overline{n} =_β \overline{\mathsf{ev}(\mathtt{f}, n)} = \overline{f(n)}.
:::

:::Remark
While our system appears to be restricted to
elementary recursive functions,
it turns out that there is another dependently
typed system in which functions on natural numbers are
*primitive* recursive [@buchholtz2024primitive].
Interestingly enough, part of their solution is in
forbidding function types from appearing in the
lowest universe, just like in %PTSrk{}.
:::

There is also a boundary result stemming from
information-theoretical considerations.

:::Theorem {#rep-bound}
If a function $f : ℕ^k → ℕ$ is representable in
%PTSrk{}, it is asymptotically bounded above by an
elementary recursive function of the size of the input
encoding.
:::

:::Proof
Since it is representable, there are
$N_I : \rk{r}$, $N_O : \rk{s}$ and
$F : N_I → … → N_I → N_O$ such that $F$ represents
$f$.
Therefore,

%mathpar
  ∀ (n_1, …, n_k), \quad
  F\,[n_1]^I\,…\,[n_k]^I =_β [f(n_1, …, n_k)]^O

Since terms are definitionally equal,
they must have the same normal form which is then
equal to $[f(n_1, …, n_k)]^O$ as a word.
Looking at the proof of Theorem <#bounded-sn>,
its length is then bounded above by
$\mathcal{O}(c ↑ … ↑ c ↑ i(n_1, …, n_k))$ where
$i(n_1, …, n_k)$ is the length of
$F\,[n_1]^I\,…\,[n_k]^I$ as a word.
The height of the tower of exponents depends not on
the $n_i$-s, but on the $d$-bound which is equal
to $r + 1$ so this bound is %textsc{ELEMENTARY}.

By combinatorial considerations,
length of $[x]^O$ is $Ω(\log{x})$ so

%multline*
  f(n_1, …, n_k) = \mathcal{O}(e ↑ \left(c
  ↑ … ↑ c ↑ i(n_1, …,
           n_k)\right)) = \\ = \mathcal{O}(e ↑ c ↑ … ↑ c
  ↑ i(n_1, …, n_k)),

thus bounded by an elementary recursive function
of the size of input encoding.
:::
In a particular case where elementary recursive
encoding is used, function representable in
%PTSrk{} is bounded by elementary recursive
function.

:::Corollary
Functions representable in %PTSrk{} with
elementary recursive bound on input encoding size
are precisely elementary recursive functions.
:::

:::Proof
Since functions computable under elementary
recursive time bound are elementary recursive [@CookKapron2017PrimitiveRecursive],
this follows from Theorems <#non-uniform-rep>,
<#rep-bound>.
:::

Two questions remain open: a) whether there are
non-elementary functions which are non-uniformly
computable in %PTSrk{} because of non-elementary
input encoding and b) whether there are any
functions which are not in $\mathcal{E}^2$ and are
uniformly computable in %PTSrk{}.

# Set-theoretical model of %PTSrk{} {#rk-model}

Now, studying computational properties of a
calculus is great, but logic is only useful when
it can tell right from wrong.
In this section, we are going to provide
nontrivial set-theoretical semantics for %PTSrk{}
with respect to which it is sound,
and where $⟦Π (X : \rk{j}).X⟧$ is going to be an
empty set, as expected.

However, we are going to need an unusual assumption:

:::Theorem
%PTSrk{} has a sound model in ZF+GCH. %formalized{rocq-ai:ZFC_model/ModelSoundness.v\#L133}
:::

:::Proof
Produced in discussion with GPT-6 Astra.

Take distinct codes $c_κ$ for every cardinal $κ$
below $ℵ_ω$ and distinct additional codes $s_j$
for $j ∈ ℕ$.
Define

%mathpar
  U_i = \{c_κ : κ ≤ ℵ_i\} ∪ \{s_j : j ∈ ℕ\},
  \qquad
  \operatorname{El}(c_κ ) = κ, \quad
  \operatorname{El}(s_j) = U_j.

That is,

%align*
  U_0 & = \{c_0, c_1, …, c_ω = c_{ℵ_0},
                  s_0, s_1, …\}, \\
          U_1 & = \{c_0, c_1, …, c_{ℵ_0},
                  c_{ℵ_1}, s_0, s_1, …\}, \\
          …\\
          U_k & = \{
              c_0, c_1, …, c_{ℵ_0},
              c_{ℵ_1}, …, c_{ℵ_k},
              s_0, s_1, …\}\\
          …

(Note that $x ∈ U_i ⇒ |\mathrm{El}(x)| ≤ ℵ_i$
since all $U_i$ are countable.)

Interpretation then goes as follows:

- Interpret $\rk{i}$ as $s_i$.
  The sort axiom holds because
  $s_i ∈ U_0 = \operatorname{El}(s_0)$.

- If $a ∈ U_m$ and $b(x) ∈ U_n$ for every
  $x ∈ \operatorname{El}(a)$, the ordinary
  dependent function set has cardinality at most

  \[
    \left|∏_{x ∈ \operatorname{El}(a)}
    \operatorname{El}(b(x))\right| ≤ ℵ_n^{ℵ_m} ≤ ℵ_{\max(m + 1, n)}.
  \]

  (The last inequality follows from GCH since
  alephs have finite indices.)

  Thus, for all such $a, b$ we can choose a
  cardinal $κ_{Π (x : a).b(x)}$ (whose
  code lies in $U_{\max(m + 1, n)}$) along with
  a bijection

  \[
    f : \left(∏_{x ∈ \operatorname{El}(a)}
    \operatorname{El}(b(x))
    \right) ≃ κ_{Π (x : a).b(x)}
  \]

  The eligible pairs $(κ, f)$ form a
  set, so ordinary choice (which holds in
  ZF+GCH [@sierpinski1947hypothese; @specker1954verallgemeinerte]) suffices.
  Make each choice
  once for each $(a, b)$, independently of a
  typing derivation or rank bound.

- Abstraction and application are modelled by $f$
  and $f^{-1}$ respectively, as defined above.

- Contexts are interpreted by valuations.

It is easy to check that the typing rules are
respected by this interpretation,
making it a sound model.
:::
A few corollaries:
- Indeed, the product interpreting
  $Π (Z : \rk{0}).Z$ ranges over $U_0$,
  which contains $c_0$, whose decoding is $∅$.
  That product is therefore empty.
- A system with rank cumulativity
  $t : \mathrm{rk}_j ⇒ t : \mathrm{rk}_{j + 1}$ is
  sound with respect to this model;
  normalization of bounded terms remains
  elementary as in Theorem <#bounded-sn>.

Note that usage of GCH is unavoidable in this
model to validate the rules.
We could try building a model in ZFC using
$\beth$-s instead of $ℵ$-s, but that would break
an assumption that $U_i$ are countable,
which would falsify axioms $\rk{n} : \rk{0}$.
That is not a huge problem, however, since ZF+GCH is
interpretable in ZF [@godel1938consistency; @godel1940consistency]
(in contrast to ZF + inaccessible cardinals which
is not [@kanamori2003higher; @kunen1980set; @godel1931formal]).

# %PTSrk{} as a proof system {#proof-system}

To demonstrate the ability of %PTSrk{} to embed
nontrivial logical theories, we construct a fixed
interpretation of Tarski–Mostowski–Robinson
arithmetic $\R$ in %PTSrk{}: Church numerals and
certified finite observations first interpret its
relational form $\Rt$; definable graphs yield
$\R$.
We do so because interpretation of $\Rt$ is much
simpler than direct interpretation of $\R$ and
already exposes all the main ideas.
Direct translation is available in supplementary
formalized materials %formalized{rocq-ai:Interpretation_of_R/functional/DirectRTheorem7.v\#L75}.

The following proof was obtained in a discussion
with GPT-6 Astra.

## Classical arithmetic and the target calculus

Let $\Rt$ be the relational presentation of
Tarski–Mostowski–Robinson arithmetic in the
first-order language with equality
$\{0, S, =, ≤, A, M\}$, with ternary predicates
$A, M$ and numerals $\bar n = S^n0$.
Its axioms are the universal closures of the
following schemes, indexed by external $m, n ∈ ℕ$:

%align
  A(\bar m, \bar n, x) & ↔ \overline{m + n} = x,
  \tag{$\widetilde R1$}\label{rtilde:add}\\
  M(\bar m, \bar n, x) & ↔ \overline{mn} = x,
  \tag{$\widetilde R2$}\\
  \bar m & ≠ \bar n \quad (m ≠ n),
  \tag{$\widetilde R3$}\\
  x ≤ \bar n & ↔ ⋁_{i = 0}^{n}x = \bar i,
  \tag{$\widetilde R4$}\label{rtilde:bound}\\
  x ≤ \bar n & ∨ \bar n ≤ x.
  \tag{$\widetilde R5$}\label{rtilde:compare}

Here $m + n$ and $mn$ are calculated externally:
the language has no binary arithmetic function
symbols.
Section <#sec:nonrelational> relates this theory
to the original, functional arithmetic $\R$ [@tarski1953undecidable].

We now begin constructing the model inside
%PTSrk{} and proving that the axioms hold in our
model.
First, here are some preliminaries on the
interpretation of first-order logic formulas:

%mathpar
  ⊥ ≜ (Π (Z : \rk{0}).Z) : \rk{1}, \qquad
  ⊤ ≜ \rk{0} → \rk{0} : \rk{1}, \qquad
  \mathsf{id} ≜ λ (Z : \rk{0}).Z : ⊤.

Use the fixed-answer proof types

%gather*
  \nc P ≜ P → ⊥, \qquad
      P \ac Q ≜ (P → Q → ⊥) → ⊥, \\
      ⋁_c(P_0, …, P_k) ≜ (P_0 → ⊥) → ⋯ → (P_k → ⊥) → ⊥, \\
      P \iffc Q ≜ (P → Q) \ac (Q → P).

A type $P$ is *stable* when $\nc\nc P → P$ is
inhabited.
Negative types, dependent functions with stable
codomain, and these encoded conjunctions and
disjunctions are stable.
They support the classical logical rules on stable
proof types; case elimination is into stable
targets.
In particular, $P \oc \nc P$ has the proof
$λ (k : P → ⊥)(l : (P → ⊥) → ⊥).\,l\,k$.
We never eliminate $⊥$ into an arbitrary positive
type.

For a carrier $D$, a domain predicate $\Good$,
and fixed stable atomic interpretations,
translate quantifiers by

%align*
  (∀ x\,φ)^∗ & ≜ Π (x : D).\,\Good(x) → φ^∗ (x), \\
      (∃ x\,φ)^∗ & ≜ (Π (x : D).\,\Good(x) → φ^∗(x) → ⊥) → ⊥.

All translated formulas are stable.
Once domain closure and atomic congruence are
proved, the logical rules and equality
substitution extend structurally to formulas on
admitted arguments.
All interpretations below use this one logical
encoding.

To lift propositions $P : \rk{a}$ to
arbitrary higher ranks $k ≥ 2$,
let

%mathpar
  L_k(P) ≜ (Π (T : \rk{k - 2}).T → T) → P : \rk{k}.

Abstraction gives a map from $P$ to $L_k(P)$;
application on the term $(λ \,T\,x.x)$ gives a map in
the opposite direction.

## Unary domain and equality

Use Church numerals with addition and
multiplication as a carrier: $D ≜ ℕ_1$.
As a reminder, here are two only *constructors* of
Church numerals, $\z$ and $\suc$:

%mathpar
  \z ≜ λ (X : \rk{1})(f : X → X)(x : X).x \quad
  \suc ≜ λ (n : ℕ_1)\,X\,f\,x.f\,(n\,X\,f\,x).

Even if the theory we interpret and the host type
system are not strong enough to support induction
on Church numerals, we still need to prove some
properties by induction.
For this, we introduce a special helper
proposition as a domain predicate $\Good$.
It would state for a given $d : V$ that each
inductive property $Q$ is satisfied for $d$:

%multline*
  \Good(d) ≜ \bigl(Π (Q : D → \rk{10}).\;
      Q(\z) → (Π (u : D).\,Q(u) → Q(\suc u)) → Q(d)\bigr)\\ : \rk{12}.

Zero and successor preserve it by abstraction.
For example,

%mathpar
  \mathsf{good}_0 ≜ λ Q\,q_0\,q_s.
      \,q_0, \qquad
  \mathsf{good}_s\,d\,g ≜ λ Q\,q_0\,q_s.\,
      q_s\,d\,(g\,Q\,q_0\,q_s).

Define equality by negative tests:

%mathpar
  E(d, e) ≜ Π (P : D → \rk{1}).\;
  (\nc P(d)) \iffc (\nc P(e)) \quad : \rk{5}.

It is reflexive, symmetric and transitive;
every $F : D → D$ respects it.
Equality substitution for $E$ itself follows from
symmetry and transitivity.
Define the arithmetic atoms by

%mathpar
  A^∗ (d, e, f) ≜ E(\Add\,d\,e, f), \qquad
  M^∗ (d, e, f) ≜ E(\Mul\,d\,e, f).

They respect $E$ in all arguments;
axioms $\widetilde R1$ and $\widetilde R2$ hold
definitionally.

## Order from successor-closed unary tests

:::Definition
A *certified unary algebra* is the fixed telescope

%align*
  \mathcal U = {} & (X : \rk{1})\,(z : X)\,
  (s : X → X)\,(δ : X → \rk{4})\\
  & (d_0 : δ z)\,
  (d_s : Π (u : X).\,δ u → δ (su)).
:::

Basically, this telescope contains all the
necessary information to “reinterpret” Church
numerals into new carrier $X$ with a domain
predicate $δ$.

Writing $\mathcal U → B$ abbreviates the
successive dependent arrows for these binders;
it does not introduce a record type.
Put $\ev_{\mathcal U}(d) = d\,X\,s\,z$ and
$G(d) ≜ (\mathcal U → δ (\ev_{\mathcal U}(d))) : \rk{6}$.
For $P : X → \rk{1}$, write

%mathpar
  \Up_{\mathcal U}(P) ≜ Π (u : X).\,δ u → \nc P(u) → \nc P(su).

This proposition asserts that a negation of $P$ is
upwards closed for ‘good’ values in the new domain $X$.
Define

%align*
  O(d, e) ≜ \bigl(
  \mathcal U → Π (P : X → \rk{1}).\; &
  \Up_{\mathcal U}(P) → {}\\
  & \nc P(\ev_{\mathcal U}(d)) → \nc P(\ev_{\mathcal U}(e))\bigr) : \rk{6}.

With this definition, preorder laws and equality
compatibility can be derived.

:::Lemma[\formalized{rocq-ai:Interpretation_of_R/relational/DirectRUnary.v}] {#r-common}
The following holds:

- $G(\z)$, $G(d) → G(\suc d)$ and
  $\Good(d) → G(d)$ are inhabited.

- $O$ is reflexive and transitive.

- $E(d, e) → O(d, e)$ and $G(d) → O(d, \suc d)$
  are inhabited.

- Arguments of $O$ respect $E$.
:::

:::Proof
Proof terms for most propositions are constructed
in an obvious way.
The last proposition holds because $O$ is
transitive and $E$ is symmetric.
:::

Consequently, for each external $m ≤ n$,
there is a closed proof of
$O(\overline{m}, \overline{n})$, built by a finite
succession of steps.

## Finite separation {#finite-separation}

In this section, we prove that $\widetilde R3$ also
holds in our interpretation of $\Rt$.

:::Lemma[\formalized{rocq-ai:Interpretation_of_R/relational/DirectRSeparation.v\#L143}]
For each $n ∈ ℕ$, there is a proof of
$O(\overline{n + 1}, \overline{n}) → ⊥$.
:::

:::Proof
By contradiction.
Suppose $H : O(\overline{n + 1}, \overline{n})$;
put $c = n + 1$, $q = c + 1$,
and remember already familiar type of selectors

%mathpar
  F_q ≜ \underbrace{\rk{1} → ⋯ → \rk{1}}_{q\ \mathrm{arguments}} → \rk{1}, \qquad
  e_i ≜ λ Z_0… Z_{q - 1}.Z_i.

We are going to use them as a “model” to interpret
an assumed $O$ and get a contradiction:
set $X = F_q$, $z = e_0$, and use “saturated
addition” as a successor:

%mathpar
  s(u) ≜ λ Z_0… Z_{q - 1}.\
  u\,Z_{\min(0 + 1, c)}⋯ Z_{\min(c + 1, c)}.

Thus, $s(e_i) =_β e_{\min(i + 1, c)}$.
Define the domain predicate as

%mathpar
  δ (u) ≜ \bigl(Π (Q : F_q → \rk{3}).\;
  Q(e_0) → ⋯ → Q(e_c) → Q(u)\bigr) : \rk{4}.

Selector projection gives $d_0$.
To construct $d_s$, open the output $δ (su)$ at
$Q$ and its branch proofs, and specialize the
input certificate to $v ↦ Q(sv)$.
Its branches are the supplied proofs for the
appropriate successor selectors.

Use the unary test

%mathpar
  P(u) ≜ u\,\underbrace{⊤⋯⊤}_{c\ \mathrm{entries}}\,⊥.

To prove $\Up(P)$, eliminate $δ (u)$ into the
rank-three motive $v ↦ \nc P(v) → \nc P(sv)$.
At $e_c$ both tests are $⊥$ and the implication is
identity.
At $e_i$ with $i < c$, there is a premise $⊤ → ⊥$.

Specializing $H : O(\overline{c}, \overline{n})$
with these data yields $(⊥ → ⊥) → (⊤ → ⊥)$.
:::

:::Lemma[\formalized{rocq-ai:Interpretation_of_R/relational/DirectRSeparation.v\#L193}] {#sep}
For each $m > n$, there is a proof of
$O(\overline{m}, \overline{n}) → ⊥$.
:::

:::Proof
Suppose $O(\overline{m}, \overline{n})$ is
inhabited; since $m > n$,
$O(\overline{n + 1}, \overline{m})$ is also
inhabited.
By transitivity of $O$ it means that
$O(\overline{n + 1}, \overline{n})$ which we know
is absurd.
:::

Thus, $\widetilde R3$ also holds: if $m ≠ n$,
then either $m > n$ or $m < n$ (these are external
natural numbers).
Assuming $E(\overline{m}, \overline{n})$,
we obtain a paradoxical $O$ and derive a
contradiction by Lemma <#sep>.

## Bounded cases {#bounded-cases}

Here, we prove a crucial lemma from which
remaining axioms follow.
Set

%mathpar
  \Case_n(d) ≜ ⋁_c
  \bigl(E(d, \overline{0}), …, E(d, \overline{n}),
  O(\overline{n + 1}, d)\bigr) : \rk{8}.

:::Lemma[\formalized{rocq-ai:Interpretation_of_R/relational/DirectRUnary.v\#L808}] {#case}
$Π (d : D).\;\Good(d) → \Case_n(d)$ is inhabited
for external $n$.
:::

:::Proof
We are going to use the following positive
invariant for $Q$:

%mathpar
  \Inv_n(d) ≜ Π (Z : \rk{8}).\;(G(d) → \Case_n(d) → Z) → Z \quad : \rk{10}.

At zero, supply $G(\z)$ and the
$E(\z, \overline{0})$ branch.
At successor, open the target continuation and
eliminate the child invariant at that same $Z$,
obtaining $G(d)$ and $\Case_n(d)$.
Obtain $G(\suc d)$ independently.
To construct the new negative case proof,
all branch continuations return $⊥$.
An $i < n$ branch becomes $i + 1$ by successor
congruence.
The $i = n$ branch gives
$E(\suc d, \overline{n + 1})$,
hence $O(\overline{n + 1}, \suc d)$.
The large branch is composed with $O(d, \suc d)$.
This proves successor closure.

Instantiate $\Good(d)$ with $\Inv_n$.
Project its case component by choosing
$Z = \Case_n(d)$ and continuation $λ g\,c.c$.
No negative disjunction is eliminated into
positive $G$-evidence.
:::

## The relational interpretation

:::Theorem[\formalized{rocq-ai:Interpretation_of_R/relational/DirectRTheorem4.v\#L57}] {#main}
^[The linked theorem gives original-PTS
inhabitants for all axiom instances and all source
derivations.
The generic constructive witness search turns each
such existence proof into a proof term together
with its typing derivation.]
The carrier $D$, admitted domain $\Good$,
closed operations $\z, \suc$,
and relations $E, O, A^∗, M^∗$ interpret the
classical relational theory $\Rt$ in %PTSrk{}.
Every translated axiom has an effectively
constructible closed proof term.
The interpretation is closed under classical
first-order derivations: every $\Rt$-derivation of
a sentence effectively yields a closed %PTSrk{}
proof term of its translation.
True ground arithmetic atoms have proofs,
and false ground atoms have internal refutations.
:::

:::Proof
Domain closure, atomic stability,
and congruence have been established;
hence source deductions translate to PTS proofs.
We know that $\widetilde R1$-$\widetilde R3$ hold;
let us check the remaining axioms.

For bounded order ($\widetilde R4$),
apply Lemma <#case>.
Under $O(d, \overline{n})$, a case involving
$O(\overline{n + 1}, d)$ contradicts
$\nc O(\overline{n + 1}, \overline{n})$ by
transitivity.
Conversely, a $E(d, \overline{i})$ branch with
$i ≤ n$ yields $O(d, \overline{n})$ via $E ⇒ O$
and ground comparison.

For comparability ($\widetilde R5$),
the finite branches of Lemma <#case> give
$O(d, \overline{n})$ and the large branch gives
$O(\overline{n}, d)$.

The ground proof/refutation claims follow from the
arithmetic beta equations, reflexivity,
ground order, and Lemma <#sep>.
:::

:::Remark
Constructions from Sections <#finite-separation>
and <#bounded-cases> generalise to arbitrary
Church-encoded inductive types providing
constructor disjointness and injectivity,
as well as finite case analysis for subtrees of
fixed constructor trees.
These results can be seen as internal inversion
lemmas.
:::

While this interpretation is full, it, however,
is not faithful on arbitrary first-order
sentences.

:::Remark[Failure of faithfulness \formalized{rocq-ai:Interpretation_of_R/relational/DirectRTheorem4.v\#L47}]
The interpretation of $\Rt$ is not faithful.
:::

:::Proof
The translation of $ρ ≜ ∀ x\,(x ≤ x)$ is provable
in %PTSrk{}, but $\Rt ⊬ ρ$.

We already know that $O$ is reflexive (Lemma <#r-common>).
But why is reflexivity not derivable in
$\widetilde R$?
It follows from model-theoretic considerations:
there is a set-theoretic countermodel.

Take $N^+ = ℕ ⊔ \{ω\}$, where $ω$ is one new
element.
Interpret $0$ and successor as usual on $ℕ$,
put $Sω = ω$, and use actual equality.
For $x, y ∈ ℕ$, let $A(x, y, z)$ hold exactly when
$z = x + y$, and let $M(x, y, z)$ hold exactly
when $z = xy$.
When at least one of $x, y$ is $ω$,
let each predicate hold exactly when $z = 0$.
Define

%mathpar
  x ≤^+ y \quad ⟺ \quad
  \bigl(x, y ∈ ℕ\ \text{and}
  \ x ≤_{ℕ} y\bigr)
  \ \text{or}\ \bigl(x ∈ ℕ\ \text{and}
  \ y = ω \bigr).

The ground graph schemes hold for every third
argument, and distinct numerals are unequal.
For every standard $n$, the downset of $n$ is
exactly $\{0, …, n\}$, proving both directions of
(<#rtilde:bound>).
Every standard number is comparable with $n$,
and $n ≤^+ ω$, proving (<#rtilde:compare>).
Hence, this is a model of $\Rt$.

However, $ω \not ≤^+ ω$, so $ρ$ is false.
:::

Actually, it is possible to prove inside %PTSrk{}
that $\Good$ Church numerals form a *strict ordered
commutative unital semiring*.
A generated proof is available in
supplementary materials.


## Relation to the original presentation of %texorpdfstring{$\R$}{R} {#sec:nonrelational}
The original Tarski–Mostowski–Robinson
arithmetic $\R$ [@tarski1953undecidable] uses
$\{0, S,+, ⋅, ≤ \}$ and the following schemes,
again with universal closure:

%align
  \bar m + \bar n & = \overline{m + n}, & & \tag{R1}\\
  \bar m ⋅ \bar n & = \overline{mn}, & & \tag{R2}\\
  \bar m & ≠ \bar n \quad (m ≠ n), & & \tag{R3}\\
  x ≤ \bar n & ⟶
      ⋁_{i = 0}^{n}x = \bar i,
      & & \tag{R4}\label{r:bound}\\
  x ≤ \bar n & ∨ \bar n ≤ x.
      & & \tag{R5}\label{r:compare}

Requiring $A, M$ in $\Rt$ to be total,
single-valued graphs makes them interchangeable
with the function symbols $+, ⋅$.
This functional strengthening is definitionally
equivalent to $\R$; even without that
strengthening, the two theories have the same
interpretability degree.

:::Proposition
The classical theories $\Rt$ and $\R$ are mutually
interpretable ^[Mutual interpretability does not
imply bi-interpretability.
In particular, the totalized graphs need not
coincide with the original relations $A, M$ on
arbitrary arguments.].
:::

:::Proof
In $\R$, interpret $A(x, y, z)$ by $x + y = z$ and
$M(x, y, z)$ by $xy = z$.
The ground arithmetic equations give the two graph
schemes.
For $i < j$, an assumption $\bar j ≤ \bar i$
contradicts (<#r:bound>) and numeral separation;
(<#r:compare>) therefore yields $\bar i ≤ \bar j$.
The case $i = j$ follows from (<#r:compare>)
directly.
Thus, $\R$ proves all true ground comparisons,
and equality substitution supplies the reverse
implication in (<#rtilde:bound>).
The remaining relational schemes are immediate.

For the other direction define, in $\Rt$,

%align*
  U_A(x, y, z) & ≜ A(x, y, z) ∧ ∀ w\bigl(A(x, y, w) → w = z\bigr), \\
          F_A(x, y, z) & ≜ U_A(x, y, z) ∨ \bigl(z = 0 ∧ ¬ ∃ v\,U_A(x, y, v)\bigr).

For fixed $x, y$, either a unique $A$-output
exists or it does not.
In the first case $F_A$ selects that output;
in the second it selects zero.
Hence, $F_A$ is provably total and single-valued.
At standard inputs, (<#rtilde:add>) forces its
value to be $\overline{m + n}$.
Define $F_M$ identically with $M$ in place of $A$;
its standard value is $\overline{mn}$.
These graphs interpret $+, ⋅$;
keep $0, S, =, ≤$ unchanged.
All five schemes of $\R$ follow.

Composing this graph interpretation with Theorem <#main>
gives the claimed interpretation in %PTSrk{}.
The composition uses the encoded classical
quantifiers to express values and uniqueness;
it does not assert $\Good$-closure of the raw
Church operators $\Add, \Mul$.
:::

:::Corollary
%PTSrk{} interprets $\R$ by definable arithmetic
graphs.
:::

# Conclusion

To sum it all up, what started as a funny idea
(what if there actually was a sound pure type
system with $\mathsf{type} : \mathsf{type}$) led
us to stumble on a truly unique dependently typed
pure type system, %PTSrk{}.

First of all, its normalization proof is
suspiciously simple and reminds of classic proofs
of normalization for STLC where,
instead of an established technique by
reducibility candidates, one patiently tracks
complexities of arrow types and makes sure that
they become simpler after sequences of reductions
(Section <#rk-sn>).
To make sure that this approach works for our
system, we formalized our proof in Rocq proof
assistant (Section <#formalization>) and in doing
so prepared the groundwork for future verified
interactions with %PTSrk{}.

After noticing that normalization proof yields a
primitive recursive normalization procedure,
we found tighter bounds and found out that it has
the same complexity as normalization of STLC,
which was previously unheard of for a dependently
typed system (Section <#complexity>).
In contrast with existing pure type systems,
our system is able to represent elementary
recursive functions (both in the sense that it is
stronger than STLC and weaker than more classical
dependently typed systems) and seems intimately
connected to it (Section <#computational>).
To check if constructions inside %PTSrk{} even
make sense, we constructed a set-theoretical model
for it which assumes a Generalized Continuum
Hypothesis to strictify the hierarchy of
$ℵ$-s which act as interpretations of our function
types (Section <#rk-model>).
Finally, we provide a lower bound on
proof-theoretical strength of %PTSrk{} by
providing a full interpretation of a relational
version of Tarski–Mostowski–Robinson arithmetic
$R$ inside it (Section <#proof-system>);
we did it partly because we suspect our system to
be able to prove its own consistency in the sense
Pakhomov finitistic set theory does [@Pakhomov2019AWeakSet],
therefore, interpretation of $R$ is the natural
first step.

%printbibliography{}
